% Part: model-theory
% Chapter: interpolation
% Section: introduction

\documentclass[../../../include/open-logic-section]{subfiles}

\begin{document}

\olfileid{mod}{int}{int}

\olsection{Pambuka}

Teorema interpolasi iku asil mangkene: Upamane
$\Entails !A \lif !B$. Banjur ana !!a{sentence} $!C$ kang
$\Entails !A \lif !C$ lan $\Entails !C \lif !B$. Kajaba iku, saben
!!{constant}, !!{function}, lan !!{predicate} (saliyane $\eq$) ing
$!C$ uga ana ing $!A$ lan~$!B$. !!{sentence} $!C$ iki diarani
\emph{interpolan} kanggo $!A$ lan~$!B$.

Teorema interpolasi narik kawigaten kanthi dhewe, nanging
wigati utamane amarga bisa digunakake kanggo mbuktekake asil
ngenani apa kang bisa ditemtokake sajrone sawijining teori, lan
syarat-syarat supaya gabungan rong teori kang saben-saben
konsisten uga konsisten. Asil kapisan dikenal minangka
teorema Beth ngenani apa kang bisa ditemtokake; asil kapindho,
teorema konsistensi bebarengan Robinson.

\end{document}
